% =====================================================================
%  CHƯƠNG 1 — MỞ ĐẦU
% =====================================================================
\chapter{Mở đầu}

\section{Đặt vấn đề và tính cấp thiết}

Ngành bán lẻ vận hành trên ba đặc điểm tài chính đi kèm với nhau: biên lợi nhuận ròng mỏng, hàng
tồn kho và vốn lưu động chiếm tỉ trọng lớn, và chi phí thuê mặt bằng cố định. Trong bộ dữ liệu của
khóa luận, các con số này quan sát được trực tiếp: tỷ số nợ trên tổng tài sản (\texttt{debt\_to\_assets})
có trung vị 0,72 — hơn một nửa số quý khảo sát có nợ phải trả chiếm trên 70\% tổng tài sản; tỷ số
hàng tồn kho trên doanh thu (\texttt{inventory\_to\_sales}) có trung vị 0,64. Khi doanh thu đi ngang
hoặc giảm, đòn bẩy cao khuếch đại rủi ro rất nhanh: lợi nhuận hoạt động sụt nhẹ cũng đủ đẩy hệ số
khả năng thanh toán xuống dưới ngưỡng an toàn.

Đối với chuỗi bán lẻ, hậu quả của việc phát hiện muộn là không đối xứng. Chậm một đến hai quý có
thể đồng nghĩa mất khả năng đàm phán với nhà cung cấp, phải giảm giá để giải phóng hàng tồn, hoặc
huy động vốn ở điều kiện bất lợi. Với phía cho vay và đầu tư, một ca bỏ sót (false negative) có chi
phí lớn hơn nhiều so với một báo động giả (false positive). Khóa luận định lượng bất đối xứng này
bằng giả định chi phí $C_{FN} = 5$ và $C_{FP} = 1$, và dùng chính giả định đó để chọn ngưỡng quyết
định thay vì mặc định lấy 0,5.

\section{Mục đích nghiên cứu}

Khóa luận hướng tới ba sản phẩm cụ thể thay vì chỉ báo cáo một điểm số mô hình:

\begin{enumerate}
\item Xây dựng pipeline dữ liệu tái lập được: tải \texttt{companyfacts} từ SEC, chuẩn hoá về chuỗi
      16 chỉ tiêu theo quý, gán nhãn và chia tập có dải purge chống rò rỉ. Mọi bước kiểm chứng được
      bằng SHA-256 và bằng đối chiếu ngược về nguồn công bố.
\item Xây dựng bộ đặc trưng tài chính có ý nghĩa kinh tế: tỷ số thanh khoản, đòn bẩy, sinh lời,
      hiệu quả hoạt động, tăng trưởng so với cùng kỳ và đặc trưng đường đi theo cửa sổ quý.
\item Huấn luyện và so sánh các họ mô hình phân loại, chọn cấu hình tốt nhất theo giao thức đánh
      giá chống thổi phồng (cross-company), và báo cáo trung thực cả những trường hợp mô hình không
      hơn baseline.
\end{enumerate}

\section{Đối tượng và phạm vi nghiên cứu}

Đối tượng nghiên cứu là các doanh nghiệp bán lẻ niêm yết tại Mỹ có báo cáo tài chính quý theo mẫu
10-Q và báo cáo thường niên theo mẫu 10-K. Bộ dữ liệu được giới hạn ở 8 công ty (WMT, HD, LOW,
ROST, DG, ORLY, DKS, FIVE) với 332 quý tài chính, giai đoạn FY2015--FY2026. Tiêu chí chọn công ty
gồm: (i) có tối thiểu 32 quý (8 năm) dữ liệu để đủ dựng train/validation/test theo chính sách chia
tập; (ii) đủ 16 chỉ tiêu chuẩn; (iii) không đổi kỳ tài chính giữa chừng.

\subsection{Vì sao dùng dữ liệu XBRL của SEC}

Thẻ XBRL (taxonomy) gắn cho từng dòng số liệu cho phép bóc tách theo \texttt{tag}, theo kỳ
(\texttt{start}/\texttt{end}) và theo ngày công bố (\texttt{filed}) một cách xác định. Nhờ đó việc
trích xuất không còn phụ thuộc vào việc đọc hiểu bố cục báo cáo, và mọi con số đều truy vết được
về hồ sơ gốc.

\subsection{Giới hạn phạm vi}

Khóa luận không dự báo phá sản. Kiểm tra hồ sơ 8-K cho thấy trong 8/8 công ty, số sự kiện phá sản
(8-K item 1.03) bằng 0, nên bài toán được định vị là dự báo suy giảm tài chính. Ngoài ra, vì chỉ có
8 thực thể, mọi kết luận đều phải đi kèm khoảng tin cậy và không suy rộng ra toàn ngành.

\section{Phương pháp tiếp cận}

Quy trình được tổ chức theo trình tự dữ liệu $\rightarrow$ đặc trưng $\rightarrow$ mô hình
$\rightarrow$ đánh giá. Điểm kiểm soát quan trọng nhất nằm ở khâu đánh giá, xuất phát từ một quan
sát thực nghiệm: nhãn hầu như là thuộc tính của từng công ty (HD, LOW, WMT có 100\% quý mang nhãn 1;
ROST chỉ 4,7\%). Nếu chia tập ngẫu nhiên theo mẫu, cùng một công ty sẽ xuất hiện ở cả train lẫn test
và mô hình chỉ cần nhớ "mặt" công ty là đạt điểm cao giả tạo.

Vì vậy khóa luận áp dụng đồng thời ba lớp kiểm soát:

\begin{enumerate}
\item Chia tập theo thời gian trong từng công ty (8 quý cuối là test, 4 quý trước là validation),
      kèm dải purge ở hai biên để tránh chồng lấn cửa sổ lịch sử.
\item Kiểm định chéo chia theo nhóm công ty: \texttt{GroupKFold} và \texttt{Leave-One-Company-Out}.
\item Luôn đặt cạnh một baseline không dùng mô hình — lấy tỉ lệ nhãn trung bình theo công ty trong
      train — để biết phần hiệu năng nào thực sự đến từ năng lực dự báo.
\end{enumerate}

Toàn bộ mã nguồn, artifact và kết quả trong khóa luận tái lập được bằng
\texttt{python -m scripts.run\_all}.
